\chapter{Foundation Models}

The Transformer architecture is the substrate on which modern language models---and
therefore agentic AI systems---are built. This chapter explains the modelling foundations
needed for the rest of the book: how text becomes tokens, how tokens become vectors, and
why decoder-only Transformers became the default architecture for general-purpose
assistants.

\section{Tokenisation: the First Compression Step}

Before a model can process text, that text must be converted into tokens: discrete units
such as words, subwords, bytes, or characters. Tokenisation shapes cost, context length,
and multilingual behaviour, and is therefore a consequential design decision rather than
a minor preprocessing detail.

Subword tokenisation, implemented through algorithms such as Byte-Pair Encoding, is the
approach used by most modern language models. Common words are represented as a single
token; rare words are decomposed into reusable subword pieces; and out-of-vocabulary
terms are handled gracefully by decomposing them into familiar subword units, balancing
vocabulary size against sequence length more effectively than word-level or
character-level alternatives.

For agent builders, tokenisation has direct engineering consequences: the cost and
context capacity of any operation scale with token count rather than character count,
making token-aware context engineering a production requirement.

\section{Embeddings and Representation Geometry}

After tokenisation, each token is mapped to a dense vector called an embedding. Dense
embeddings place tokens in a learned vector space where similar meanings, syntactic
roles, and usage patterns tend to occupy related regions, expressing semantic relatedness
as geometric proximity.

The same principle underlies retrieval. A document chunk and a query can both be
embedded into the same vector space, and similarity in that space selects relevant
context. Both language models and retrieval-augmented generation systems rest on the
assumption that meaning can be approximated by geometry in a learned representation
space.

\section{The Limitations of Sequential Models}

Before Transformers, recurrent models such as RNNs and LSTMs processed text
sequentially, updating a hidden state token by token. This introduced order sensitivity
but imposed two significant constraints: long-range dependencies were difficult to
maintain as information from early tokens could fade across the sequence, and training
was difficult to parallelise because later hidden states depended on earlier ones.

Agentic workloads routinely require long-range dependencies---connecting a term defined
early in a contract to its use later, or a failing test to the function responsible for
it across multiple files. The Transformer addressed both constraints simultaneously, and
its dominance in production systems reflects how consequential those constraints were in
practice.

\section{Self-Attention}

Self-attention allows each token to directly compare itself to every other token in the
sequence. For each position, the model constructs three learned projections---a query, a
key, and a value---where the query represents what the token is looking for, the key
represents what another token offers for matching, and the value represents the content
to be aggregated.

Scaled dot-product attention is defined as:

\[
\operatorname{Attention}(Q,K,V)
= \operatorname{softmax}\!\left(\frac{QK^{\top}}{\sqrt{d_k}}\right)V.
\]

The matrix $QK^{\top}$ contains pairwise relevance scores between query and key
positions; the softmax converts those scores into a probability distribution; multiplying
by $V$ produces a weighted combination of value vectors; and the division by
$\sqrt{d_k}$ stabilises the scale of dot products as dimensionality grows.

Multi-head attention repeats this operation with multiple learned projections in
parallel, allowing the model to attend simultaneously to local syntactic structure,
entity references, and long-range answer spans, producing representations that are
context-aware across multiple dimensions at once.

\section{The Original Transformer}

The original Transformer, introduced in \emph{Attention Is All You Need}, used an
encoder-decoder architecture. The encoder processed the source sequence with full
bidirectional self-attention; the decoder generated the target sequence autoregressively
through masked self-attention over prior outputs and cross-attention over encoder
representations.

The architecture established the design patterns that remain central to modern models:
attention sublayers for token-to-token interaction, position-wise feedforward sublayers
for nonlinear transformation, residual connections for gradient flow, normalisation for
training stability, positional encodings to supply order information that
permutation-invariant attention cannot recover, and causal masking to enforce
left-to-right generation. Together these choices gave every token a direct path to every
other token and expressed the computation as matrix operations that GPUs can parallelise
efficiently.

\section{From Encoder-Decoder to Decoder-Only Models}

Three Transformer families emerged from the original design. Encoder-only models such as
BERT apply bidirectional self-attention to produce strong contextual representations for
classification, extraction, and ranking. Encoder-decoder models such as T5 retain the
full sequence-to-sequence structure for tasks requiring one sequence to be transformed
into another. Decoder-only models became dominant for large generative models because
next-token prediction is simple, scalable, and supports open-ended generation.

Decoder-only models are the architecture of choice for agent systems because a
conversation, a plan, a tool call, a tool result, and a final response can all be
serialised into a single sequence that the model learns to continue in the desired
format.

\section{Positional Information}

Self-attention is permutation-invariant: the same scores result regardless of token
order unless positional information is explicitly introduced. Early Transformers added
absolute positional encodings to token embeddings; modern frontier models use Rotary
Positional Embeddings (RoPE), which encode position through the attention computation
itself and generalise more effectively to long sequences.

For agentic systems, positional design carries practical significance because context
windows spanning hundreds of thousands of tokens are now common, and the quality of
positional representations at those lengths directly affects reasoning quality.

\section{Inference: Logits, Decoding, and Structured Output}

At inference time, a decoder-only model produces a distribution over the vocabulary at
each step as a vector of logits. Greedy decoding selects the highest-probability token;
sampling introduces stochasticity controlled by temperature, top-$k$, and top-$p$
parameters.

Structured agent outputs require additional discipline. When a model must produce a
valid JSON object for a tool call, the system may rely on prompting conventions, retry
logic after schema validation failure, or constrained decoding that restricts generation
to schema-valid tokens at each step---a trade-off between flexibility, latency, and
conformance reliability.

\section{Why This Matters for Agents}

Every agentic capability is built on these modelling primitives: tool calls are generated
token sequences, retrieved documents and memory summaries are context tokens, plans are
generated text, and cost and latency scale with total token count. Understanding
Transformers at this level supports the engineering decisions that follow in every
subsequent chapter---how to manage context, when to constrain generation, why retrieval
precision matters, and how model capability translates into agent reliability.

\section{Summary}

Transformers replaced sequential recurrence with parallel self-attention, giving every
token a direct path to every other token in a form that modern hardware can execute
efficiently. Decoder-only Transformers became the default for large generative models
because next-token prediction scales cleanly and supports open-ended generation. Agents
are built by placing this generative engine inside a stateful loop with tools, memory,
context engineering, and security boundaries---the layers this book addresses in the
chapters that follow.